% !PRINTED 415
\weberpage{415}
\begin{center}
\S.\,110.

\gesperrt{The Icosahedral Equation.}
\end{center}

The result of these considerations is, in the first place, so far as it
concerns the general equations of higher than the 4th degree, a negative
one. We do not succeed, by the rational formation of resolvents, in
reducing the general equation of the 5th degree to an equation with \gesperrt{one}
parameter. If we nevertheless do not yet wish to give up this goal,
nothing remains but that we no longer regard the variables
$x_0, x_1, \ldots, x_{n-1}$ as completely independent variables, and that
we let certain equations subsist between them. We are then, to be sure,
no longer dealing with the general equation of the degree in question,
but with a more special one, in which the coefficients are not
independent of one another, and the second question presents itself, how
and by what irrational functions the general equation can be reduced to
this special one. If this succeeds by means of equations of lower degree,
e.g. by square roots, a reduction of the problem will at any rate have
been achieved.

We therefore now assume that the variables $x_0, x_1, \ldots, x_{n-1}$
are connected with one another by an arbitrary number of relations
\weq{1}{A(x_0, x_1, \ldots, x_{n-1}) = 0, \quad B(x_0, x_1, \ldots, x_{n-1}) = 0, \ldots}
in which the $A, B, \ldots$ are rational functions of the $x$, which are
so constituted in number and mutual dependence that the values of the $x$
are not numerically fixed, but are still variable. Moreover, we take as
a basis a permutation group $\mathfrak{A}$ of the
$x_0, x_1, \ldots, x_{n-1}$, and regard the functions of the $x$ that
admit the permutations of $\mathfrak{A}$, and only these, as rational,
so that $\mathfrak{A}$ is the Galois group of the equation whose roots
are the $x_0, x_1, \ldots, x_{n-1}$. \gesperrt{The relations (1) must then be so
constituted that they too admit the permutations $\mathfrak{A}$.}

Under these presuppositions we ask after the possibility of determining
two rational functions $u, z$ of $x_0, x_1, \ldots, x_{n-1}$ in such a
way that between them an equation
\weq{2}{F(u, z) = 0}
%CONT>>
